\documentclass[11pt,a4paper]{article}

\usepackage[utf8]{inputenc}
\usepackage[T1]{fontenc}
\usepackage[spanish, es-noquoting]{babel}
\usepackage{booktabs}
\usepackage{geometry}
\usepackage{hyperref}
\usepackage{xcolor}
\usepackage{fancyhdr}
\usepackage{enumitem}
\usepackage{tabularx}

\geometry{margin=2.5cm}
\setlength{\parindent}{0pt}
\setlength{\parskip}{0.4em}

\hypersetup{
    colorlinks=true,
    linkcolor=black,
    urlcolor=blue,
    citecolor=black
}

\pagestyle{fancy}
\fancyhf{}
\fancyhead[L]{\small Skalar Inc}
\fancyhead[R]{\small Revisi\'on de Datos Transaccionales -- Abril 2026}
\fancyfoot[C]{\thepage}

\newcommand{\code}[1]{\texttt{#1}}
\newcommand{\field}[1]{\texttt{#1}}

\title{
    \vspace{-2cm}
    {\LARGE Revisi\'on de la Primera Entrega de Datos \\
    Transaccionales Fudo -- Skalar} \\[0.3cm]
}
\author{Skalar Inc.}
\date{16 de abril de 2026}

\begin{document}

\maketitle
\thispagestyle{fancy}

%% ============================================================
\section{Alcance}

Este documento resume los hallazgos de la auditor\'ia de calidad realizada
por Skalar sobre la primera entrega de datos transaccionales de Fudo
(export Metabase del 16 de abril de 2026, 33\,959~filas). La revisi\'on
evalu\'o conformidad de tipos, dominios de valores, integridad de grano
y formato de entrega, seg\'un lo establecido en el \textit{Protocolo de
Comunicaci\'on de Datos Transaccionales Fudo--Skalar} (2026-04-09).

Adicionalmente, se introduce el \textbf{contrato de datos}
(\code{data\_contract.yaml}) como especificaci\'on t\'ecnica
leg\'ible por m\'aquina que acompa\~na al protocolo.

%% ============================================================
\section{Resultado de la Revisi\'on}

\begin{tabularx}{\textwidth}{l c X}
\toprule
\textbf{Verificaci\'on} & \textbf{Estado} & \textbf{Detalle} \\
\midrule
Unicidad de grano & PASS &
  Cero duplicados en (\field{client\_id}, \field{transaction\_date},
  \field{product}). \\[4pt]
Nulabilidad & PASS &
  Columnas NOT NULL sin valores nulos. \\[4pt]
Dominio de \field{currency} & PASS &
  Valores dentro de \{ARS, BRL, CLP, COP, MXN, USD\}. \\[4pt]
Dominio de \field{product} & PASS &
  Valores dentro de \{Saas, Payments, Online Ordering\}. \\[4pt]
Rango de \field{collections} & PASS &
  $\geq 0$. El 2.2\,\% de filas tiene valor cero (v\'alido). \\[4pt]
Granularidad de \field{collections} & PASS &
  Revenue reportado por producto, no agregado por cuenta. \\[4pt]
\midrule
Formato num\'erico (NF-01) & FAIL &
  Separadores de miles en campos num\'ericos; ning\'un motor
  infiere tipos correctamente. \\[4pt]
Granularidad temporal (DF-01) & FAIL &
  \field{MES} entrega periodo mensual (YYYYMM), no fecha diaria
  (YYYY-MM-DD) seg\'un \S3.1 del protocolo. \\[4pt]
Formato de fechas (DF-01) & FAIL &
  \field{FECHA\_PRIMER\_PAGO} en formato \textit{locale}
  ingl\'es, no ISO~8601. \\[4pt]
Nombre de columna (AL-01) & FAIL &
  \field{a.ACCOUNT\_ID} contiene alias SQL del query Metabase. \\[4pt]
Columnas faltantes & AUSENTE &
  \field{channel} y \field{client\_segment} no presentes en el
  export actual. \\
\bottomrule
\end{tabularx}

%% ============================================================
\section{Contrato de Datos}

El archivo \code{data\_contract.yaml} formaliza, en formato leg\'ible por
m\'aquina, las expectativas de Skalar sobre los datos que Fudo entrega.
Su prop\'osito es servir como referencia t\'ecnica para el equipo de
ingenier\'ia de Fudo durante el desarrollo del pipeline de exportaci\'on.

El contrato define:

\begin{enumerate}[leftmargin=2em]
    \item \textbf{Formato de entrega.} Reglas de codificaci\'on CSV
      (sin separadores de miles, fechas ISO~8601, UTF-8) y la
      recomendaci\'on de migrar a Parquet, formato que elimina toda
      ambig\"uedad de tipos al almacenar el esquema en metadatos del
      archivo.
    \item \textbf{Esquema can\'onico.} Las 7~columnas del protocolo con
      tipo, nulabilidad, restricciones de dominio y estado de
      conformidad actual por columna.
    \item \textbf{Expectativas de calidad.} Unicidad de grano,
      cumplimiento de nulabilidad, integridad referencial y
      completitud temporal.
    \item \textbf{Acciones priorizadas.} Lista ordenada por impacto
      (P0--P3) de los cambios que Fudo debe implementar en el export.
\end{enumerate}

Las acciones de prioridad~\textbf{P0} son:

\begin{itemize}[leftmargin=2em]
    \item Eliminar separadores de miles de todos los campos num\'ericos.
    \item Entregar \field{transaction\_date} con granularidad diaria
      (YYYY-MM-DD), utilizando el campo \field{Payment\_date} del
      sistema origen referenciado en \S3.2 del protocolo.
\end{itemize}

Skalar mantendr\'a el contrato actualizado conforme evolucione el
pipeline. Cada entrega posterior ser\'a validada autom\'aticamente
contra este contrato.

\end{document}
